\begin{textsample}{POS Dim 1 – vem\_ed\_25 – Score 28.00 – vem\_ed\_25/t136.txt}  \label{ex:f1_pos_069}
Artigo de \textbf{opinião} As WebQuests como ferramentas de sala de aula invertida Edilene Gasparini Fernandes Fatec São José do Rio Preto edilene@fatecriopreto.edu.br As WebQuests \textbf{foram} \textbf{criadas} por Bernie Dodge, em 1995, e \textbf{foram} usadas com o objetivo de integrar \textbf{disciplinas} e trazer a web para dentro das salas de aula do ensino fundamental e médio.
Elas podem \textbf{ser} consideradas como uma das ferramentas para a metodologia \textbf{ativa} conhecida como sala de aula invertida (flipped classroom), em \textbf{que} o \textbf{conteúdo} \textbf{é} apresentado fora do \textbf{ambiente} escolar e o tempo em sala de aula \textbf{é} usado pelo professor para tirar dúvidas e \textbf{promover} debates com seus estudantes.
Em 2020, as WebQuests voltaram a \textbf{ser} discutidas no Brasil por meio da pesquisa de Lilian Bacich sobre seu uso dentro do meio acadêmico (Bacich, 2020).
Motivada por esse ressurgimento, propus uma adaptação de seu uso no ensino \textbf{superior} tecnológico das Fatecs com o objetivo de melhoria do ensino-aprendizagem de língua inglesa.
Essa pesquisa \textbf{foi} financiada pelo CPS, como Regime de Jornada Integral, e envolveu alunos de três cursos diferentes da Fatec São José do Rio Preto (Tecnologia em Agronegócio, Tecnologia em Informática para Negócios e Tecnologia em Análise e Desenvolvimento de Sistemas), \textbf{que} colaboraram na criação e hospedagem das WebQuests na plataforma Google Sites.
As WebQuests \textbf{são} ferramentas \textbf{que} provocam o aluno por meio do questionamento, da curiosidade, e o motivam a buscar uma solução para o problema apresentado.
As etapas de uma WebQuest podem \textbf{variar} conforme o alvo de sua \textbf{questão}, mas, classicamente, Bernie Dodge \textbf{estabelece}: 1.
Introdução - fase em \textbf{que} se prepara o “palco” e se “fornece algumas informações de fundo”. 2.
Apresentação de uma \textbf{tarefa} “factível e interessante” \textbf{que} instigará o grupo pela curiosidade e aderência com sua \textbf{realidade} e aplicação. 3.
Indicação de endereços eletrônicos \textbf{que} \textbf{sejam} a base de partida para a aventura \textbf{que} terá outras paragens, naturalmente. 4.
Abordagem da \textbf{estrutura} de etapas na qual \textbf{será} \textbf{construída} a WebQuest: a introdução do problema ou \textbf{questão}, a \textbf{tarefa} dividida entre os membros do grupo, o \textbf{processo} de \textbf{construção} da abordagem e os recursos de \textbf{que} farão uso. \textbf{continuação} Em seguida, o projeto passa por uma avaliação interna para \textbf{que} os estudantes cheguem, juntos, à conclusão.
Depois de concluída, a WebQuest \textbf{é} também avaliada por professores ou responsáveis pelo projeto (Dodge, 1996, p. 1-2).
O resultado positivo \textbf{que} o uso das WebQuests tem proporcionado com os alunos da Fatec São José do Rio Preto ganhou também o \textbf{ambiente} dos Projetos Colaborativos Internacionais (PCIs) realizados na unidade.
Desde 2022, elas fazem \textbf{parte} dos PCIs \textbf{que} a Fatec Rio Preto realiza com países como África do Sul, Índia e Armênia.
De \textbf{início} elas \textbf{foram} empregadas como ferramentas dentro dos projetos \textbf{colaborativos}, mas a \textbf{proposta} de \textbf{tornar} a \textbf{interação} um trabalho \textbf{colaborativo} de criação de novas WebQuests \textbf{tornou} as \textbf{interações} ainda mais animadas.
Professores e alunos superaram-se em criatividade e inovação.
Em outubro de 2024, está prevista mais uma edição do PCI com alunos de Letras da Yerevan State University (Armênia), cuja \textbf{interação} lhes fornece know-how para o trabalho com as Tecnologias da Informação e da Comunicação (TICs) em \textbf{ambientes} \textbf{colaborativos}.
Yan Ribon, aluno da Fatec São José do Rio Preto envolvido com o \textbf{desenvolvimento} das WebQuests, \textbf{foi} agraciado com uma bolsa do CNPq em 2023 a \textbf{partir} de um projeto de um aplicativo para acesso mobile das WebQuests.
O aplicativo WebEnglish acabou de \textbf{ser} finalizado e está disponível, (exclusivamente no sistema operacional Android) aos colegas de trabalho e a \textbf{todos} \textbf{que} se interessarem pelo seu uso, por meio deste link: https://github.com/yanrex8/app-webquests Coloco-me à disposição para orientar quaisquer professores \textbf{que} desejem \textbf{implementar} a ferramenta em suas aulas de línguas, ou ainda, \textbf{que} pretendam trabalhar com elas em PCIs em quaisquer áreas, já \textbf{que} os assuntos \textbf{que} elas abordam \textbf{são} \textbf{variados} e envolvem diferentes áreas de \textbf{conhecimento}. \textbf{Referências} BACICH, Lilian.
WebQuest: como organizar uma atividade \textbf{significativa} de pesquisa.
Inovação na educação.
São Paulo, 22 de março de 2020.
Disponível em: https://lilianbacich.com/2020/03/22/webquest-como-organizar-uma-atividade-significativa-de-pesquisa/ Acesso em: 12 set. 24 DODGE, Bernie.
Webquest: uma técnica para aprendizagem na rede.
UFSCar, 1996.
Disponível em: https://www.dm.ufscar.br/~jpiton/downloads/artigo\_webquest\_original\_1996\_ptbr.pdf Acesso em: 12 set. 24

% matched lemmas: ambiente, ativo, colaborativo, conhecimento, construir, construção, conteúdo, continuação, criar, desenvolvimento, disciplina, estabelecer, estrutura, implementar, interação, início, opinião, parte, partir, processo, promover, proposta, que, questão, realidade, referência, ser, significativo, superior, tarefa, todo, tornar, variar
\end{textsample}
